\chapter{Lessons for Your Own Deployment}
\label{ch:lessons}

Of eleven plausible performance knobs A/B-tested on this cluster, eight were
rejected --- including several that ``obviously'' should have helped. The
single most consequential bug in the system was not a kernel fault or a
race but a configuration field that existed, defaulted sensibly, and was
never read. Results like these are why the closing chapter of this handbook
is not a victory lap: nearly every principle worth keeping was learned by
watching a reasonable assumption fail in a specific, documented way.

A handbook about one system earns its keep by what transfers out of it. This
closing chapter collects the principles the preceding twelve chapters
demonstrated in the concrete --- each stated as doctrine, grounded in the
incident that taught it, and phrased so it can be applied to a different
model, a different engine, and different hardware tomorrow.

For an inference engineer, this is the chapter meant to outlive its subject.
The specific patches are bound to one model, one engine build, and one pair
of machines; the disciplines --- pinned artifacts, fail-closed patching,
identity-keyed state, loud failure, public retraction --- are not. Each
one below carries the incident that proves what skipping it costs.

\section{Change nothing you cannot prove}

\subsection*{1. Pin everything; patch fail-closed}
The runtime image is pinned by manifest digest, the checkpoint by revision,
and every runtime modification by the SHA-256 of the bytes it expects to
find. A patch that cannot prove its target aborts the boot rather than
serving stock code under an enabled flag (\Cref{ch:patches}). The general
form: \emph{an optimization or fix that silently fails to apply is worse
than its absence}, because you will debug the system you believe you
deployed, not the one you did.

\subsection*{2. ``Defined but unenforced'' is a bug class, not a config style}
Pinning proves the bytes are present; it cannot prove a knob is connected
to anything. The starvation arrived with a clean bill of health: decode
lanes stalled behind prefill with zero preemptions and every dashboard
green. \code{max_num_partial_prefills} (\ghissue{27}) --- the cap that
should have bounded the damage --- was connected to nothing that ran. Audit the path
from every knob you rely on to the code that consumes it. The same family
produced the dead \code{IntEnum}-vs-string comparison found while auditing
\ghissue{55}.

\subsection*{3. Exactly-one spelling of ``on''}
Auditing knob-to-consumer paths is only tractable when the space of ``on''
is small. Every opt-in here is gated on the literal string \code{1}; every other value
--- \code{true}, \code{yes}, \code{2}, whitespace --- is stock. Combined with
default-off for anything unproven, this collapses the configuration state
space and makes ``what is this cluster running?'' answerable from one file
plus one rule (\Cref{ch:patches}).

\subsection*{4. Restart is not rollback}
One file plus one rule answers what the cluster is \emph{configured} to
run; what it is \emph{actually} running can still diverge.
Runtime patches live in the container's writable layer; flipping a flag off
and restarting reuses those bytes. Only recreation restores image truth ---
and the doctrine must be paired: \emph{both} ranks recreated together, never
one (\Cref{ch:operations}). Wherever you mutate a running artifact, write
down what actually reverts it, and test that path.

\section{Identity, capture, and blast radius}

The first four principles govern what you deploy; the next four govern how
it fails while serving --- and they escalate in scope, from one request's
state to the whole host.

\subsection*{5. Identity, not position}
Under continuous batching the engine condenses its running set; anything
keyed by \emph{batch row} silently reads another request's state after a
condense. DSpark's draft KV did exactly that, collapsing acceptance at
concurrency $>1$ with no crash (\Cref{ch:specdecode}). The symptom is
worth memorizing because it is so ghostly: single-stream benchmarks
perfect, multi-stream acceptance collapsing, and not one error line ---
the engine faithfully computing with another request's draft state. Every
persistent
per-request tensor needs a request-stable slot map --- and the fix should
preserve the identity-permutation fast path so the single-stream case stays
byte-exact.

\subsection*{6. Capture-time state is frozen state}
The draft-KV bug was state outliving the position that indexed it; the same
staleness has a compiled form.
Four separate patches in this system guard the same trap: a CUDA graph
captures the \emph{decisions} made during capture, not just the kernels. A
branch on live batch state, a stride from the capture-time batch, a
shortcut valid only for short contexts --- all replay wrong later, and the
result is intermittent \emph{corruption}, not a crash (\Cref{ch:bugs}).
Treat ``is this value allowed to vary after capture?'' as a review question
on every graphed path, and guard shortcuts with
\code{is_current_stream_capturing()}.

\subsection*{7. A TP pair fails as a pair}
Stale identity and frozen capture corrupt individual requests; the blast
radius widens when ranks are involved. One rank stalled --- in a JIT compile, a lost shared-memory notification, a
stochastic kernel fault --- leaves its peer blocked in a collective until a
watchdog kills both. The mitigations generalize: persist every JIT cache
across container lifecycles, warm known shapes at boot, bound every wait
(the 5-second reader recheck), and instrument for the post-mortem you will
need (the flight recorder with its archive-both-ranks rule)
(\Cref{ch:bugs,ch:operations}).

\subsection*{8. Unified memory: the host is inside the blast radius}
And the blast radius does not stop at the GPUs. On GB10, activation
workspaces, compiler caches, and desktop daemons share
physical memory with the KV pool. Two sound GPU-side tunings were rejected
purely for host-RAM cost, and the top-ranked improvement on the whole
cluster was \emph{reclaiming leaked host memory} (\Cref{ch:perf}). The
audit that ranked it found a serving node running a graphical desktop whose
polkit daemon had leaked 3.1\,GB of RSS --- bytes drawn from the same
physical pool the KV cache lives in. On any
unified-memory platform, \code{MemAvailable} and swap growth belong next to
GPU metrics in your gating checks --- and disable the desktop OOM killer.

\section{Arithmetic first, knobs second}

With the failure domains mapped, the next three principles govern how
anything about performance gets decided at all.

\subsection*{9. Derive the bound before touching a knob}
The bytes-per-step model (\Cref{ch:perf}) predicted that c=1 decode is
$\sim$80\% weight streaming --- and therefore that no configuration knob
would double it, that concurrency and acceptance were the real levers, and
that a third GPU would buy capacity rather than single-stream speed. Every
subsequent measurement, including the TP=3 natural experiment, agreed. An afternoon of
arithmetic prevented a month of misdirected tuning.

\subsection*{10. One knob per boot, gates before numbers, noise bands before verdicts}
The model says which levers can matter; measurement decides whether a given
knob moves one. The A/B campaigns changed exactly one variable per boot; verified boot gates
on both ranks before measuring; pooled eight trials because single-stream
noise was $\pm$10\,\tokspersec{}; and pre-committed verdict rules (``any drop
$\Rightarrow$ revert''). The verdicts this bought were unambiguous:
TC-decode, a knob that ``obviously'' should have helped, posted 7 of 8
trials below the baseline \emph{minimum} and was rejected on distribution
separation, not a point estimate (\Cref{ch:measurement}). Without this
discipline, the configuration would have accreted every plausible-looking
knob.

\subsection*{11. Benchmarks lie in systematic, catalogable ways}
All of that discipline still assumes the instrument itself tells the truth.
Several of the most alarming regressions this cluster ever recorded traced
to the measuring instrument, not the model --- and the pitfalls are
systematic and catalogable, not exotic. Each one produced a concrete false alarm in
\Cref{ch:measurement} before being codified into the audit suite's five
principles; steal the principles, they are engine-agnostic.

\section{Fail loud at the boundary}

Instruments are not the only boundary that can lie; the serving API is one
too, and its lies outlive the request.

\subsection*{12. Protocol errors poison agent history}
An HTTP 400 on one turn --- an image part in the wrong role, a truncated
tool call replayed --- stays in the client's transcript and fails every
subsequent request of the session. Three separate incident families
(\ghissue{55}, \ghissue{52}, \ghissue{178}) taught the same rule: signal
failure in-band the way clients actually recover (\code{finish_reason},
completed-turn semantics), and validate at the boundary the way agents
actually replay (\Cref{ch:bugs}).

\subsection*{13. Fail visible, or fail closed --- never fail plausible}
Failing loud in-band handles the errors the server notices; the harder
class is the failure it does not. The gravest failures here were the quiet
ones: HTTP 200 with zero tool calls
(\ghissue{191}), silent stream truncation with a missing
\code{finish_reason} (\ghissue{141}), acceptance collapse with perfect
uptime. The countermeasures share a shape: check the contract after
generation and prefer a 500 to a wrong 200; watch for \emph{absence}
(missing terminal fields) as the earliest signal; and put quality metrics,
not just liveness, on the dashboard.

\section{Version your beliefs}

The thirteen principles above all assume something quieter: that what the
team \emph{believes} about the system is itself maintained --- dated,
versioned, and revocable. This last discipline is the one the others
depend on, because a pinned byte proves nothing to a team that will not
un-believe a stale number.

\subsection*{14. Version your beliefs like your code}
This repository retracts in public: the ``\envvar{MAX_NUM_SEQS=32} is safe''
framing was withdrawn when repeated bursts showed stochastic per-burst
failure; an n=1 evidence matrix (\ghissue{143}) was struck; an A/B verdict
was demoted when the scheduler fix beneath it (\code{r3}) invalidated its
conditions; a patch ships with an explicit ``no rescue claim'' because the
live defect never reproduced under test. The retraction is not an apology
appended to the record --- it is architecture, carried in the same files
the configuration ships in, where the next operator cannot miss it. Dated
results, named noise bands,
best-seen-versus-norm distinctions, and evidence-status paragraphs are what
let a months-old number be trusted --- or correctly distrusted --- today
(\Cref{ch:measurement,ch:perf}).

\begin{hblesson}
If a single sentence must carry this chapter: \emph{make every claim your
system embodies --- a pinned byte, an enforced config knob, a measured
number, a safety statement --- checkable, and make every failure loud.}
The forty patches, the hash pins, the exit codes, the audit suite, and the
retractions are one idea wearing five uniforms.
\end{hblesson}

\section{A closing word}

Two desktop machines, one cable, and a public repository now serve a
million-token frontier model with vision and tools at interactive speed ---
not because any single component is miraculous, but because dozens of
unglamorous disciplines compound: a scheduler cap actually enforced, a GID
left deliberately unpinned, a cache that outlives its container, a
benchmark that refuses to warm its own prefix. That compounding is
available to any team willing to do the reading, the pinning, and the
measuring. This handbook exists so the next team starts from chapter
thirteen instead of issue \ghissue{21}.

What changes on the way out of Part~I is the unit of reuse. Almost nothing
here transfers as a patch --- the hunks are bound to one image digest, one
checkpoint revision, one attention geometry --- but the fourteen principles
travel precisely because each is a habit rather than a diff: prove the
bytes, enforce the knob, key state by identity, count the host's memory,
distrust the instrument, date every belief. Read that way, this is less a
recipe for one cluster than a checklist for the next system, including one
whose weights nobody in the lab has quantized.

\FloatBarrier
\section*{Takeaways}
\begin{itemize}
  \item Change nothing you cannot prove: artifacts pinned by digest and
    revision, patches that abort the boot rather than apply silently,
    exactly one spelling of ``on,'' and a rollback path you have actually
    tested --- recreate, because restart reuses the writable layer.
  \item Correctness in a distributed engine is identity, capture, and blast
    radius: state keyed by request rather than batch row, graphs that
    freeze nothing allowed to vary, the TP pair as one failure domain, and
    host memory counted alongside GPU memory on unified-memory silicon.
  \item Measure like an adversary and retract in public: derive the bound
    before touching a knob, change one variable per boot, prefer a loud
    failure to a plausible success --- a system's epistemics are part of
    its architecture, and a belief nobody dated cannot be trusted later.
\end{itemize}

\chaptertransition{The DeepSeek story ends here; the lab did not stop. The
same two Sparks, the same TP=2 topology over the same cable, were soon
serving a second frontier model down the opposite road --- not the vendor's
own native-precision weights but a community contributor's 4-bit EXL3
trellis quant, a DFlash2 drafter, and CUDA kernels written in-house --- so
Part~II covers only what this recipe could not teach. It opens on a
checkpoint that arrives owing proof: somebody else's arithmetic at 54\,\%
of the bytes, which must demonstrate it is not a lobotomy before it is
allowed to serve, on a stack that cannot even load its KV layout.}
